\import{algebra/ring/reducts.tex}
\import{algebra/monoid/power.tex}

\section{Natural powers in a ring}

Natural powers in a ring are exactly the generic powers in its multiplicative
monoid.  No second recursion is introduced.

\begin{definition}\label{ring_power}
    $\ringPower{R}{x}{n}=
        \monoidPower{\MultiplicativeMonoid{R}}{x}{n}$.
\end{definition}

\begin{proposition}\label{ring_power_zero}
    Let $R$ be a ring.
    Then $\ringPower{R}{x}{\emptyset}=\ringOne[R]$.
\end{proposition}
\begin{proof}
    Follows by
        \auto{ring_power,monoid_power_zero,multiplicative_monoid_fields}.
\end{proof}

\begin{proposition}\label{ring_power_type}
    Let $R$ be a ring.
    Let $x\in\carrier[R]$.
    Let $n$ be a natural number.
    Then $\ringPower{R}{x}{n}\in\carrier[R]$.
\end{proposition}
\begin{proof}
    $\MultiplicativeMonoid{R}$ is a monoid by
        \obv{multiplicative_monoid_is_exact_monoid}.
    $x\in\carrier[\MultiplicativeMonoid{R}]$ by
        \auto{multiplicative_monoid_fields}.
    $\ringPower{R}{x}{n}\in
        \carrier[\MultiplicativeMonoid{R}]$ by
        \auto{ring_power,monoid_power_type}.
    Follows by \auto{multiplicative_monoid_fields}.
\end{proof}

\begin{proposition}\label{ring_power_successor}
    Let $R$ be a ring.
    Let $x\in\carrier[R]$.
    Let $n$ be a natural number.
    Then $\ringPower{R}{x}{\suc{n}}=
        \ringMul[R](\ringPower{R}{x}{n},x)$.
\end{proposition}
\begin{proof}
    $\MultiplicativeMonoid{R}$ is a monoid by
        \obv{multiplicative_monoid_is_exact_monoid}.
    $x\in\carrier[\MultiplicativeMonoid{R}]$ by
        \auto{multiplicative_monoid_fields}.
    $\ringPower{R}{x}{\suc{n}}=
        \mul[\MultiplicativeMonoid{R}](
            \ringPower{R}{x}{n},x)$ by
        \auto{ring_power,monoid_power_successor}.
    Follows by \auto{multiplicative_monoid_fields}.
\end{proof}

\begin{proposition}\label{ring_power_one}
    Let $R$ be a ring.
    Let $x\in\carrier[R]$.
    Then $\ringPower{R}{x}{\suc{\emptyset}}=x$.
\end{proposition}
\begin{proof}
    $\MultiplicativeMonoid{R}$ is a monoid by
        \obv{multiplicative_monoid_is_exact_monoid}.
    $x\in\carrier[\MultiplicativeMonoid{R}]$ by
        \auto{multiplicative_monoid_fields}.
    Follows by \obv{ring_power,monoid_power_one}.
\end{proof}

\begin{theorem}\label{ring_power_of_negative_one}
    Let $R$ be a ring and $n$ be a natural number.
    Then $\ringPower{R}{\ringNeg[R](\ringOne[R])}{n}=\ringOne[R]$ or
        $\ringPower{R}{\ringNeg[R](\ringOne[R])}{n}
            =\ringNeg[R](\ringOne[R])$.
\end{theorem}
\begin{proof}
    $\ringOne[R]\in\carrier[R]$ by \obv{ring_one_type}.
    Hence $\ringNeg[R](\ringOne[R])\in\carrier[R]$ by \obv{ring_neg_type}.
    Let $C=\{k\in\naturals\mid
        \ringPower{R}{\ringNeg[R](\ringOne[R])}{k}
            \in\{\ringOne[R],\ringNeg[R](\ringOne[R])\}\}$.
    $\ringPower{R}{\ringNeg[R](\ringOne[R])}{\emptyset}=\ringOne[R]$ by
        \obv{ring_power_zero}.
    Hence $\emptyset\in C$ by \auto{omega_empty,upair_iff}.
    Show that for every $k\in C$ we have $\suc{k}\in C$.
    \begin{subproof}
        Fix $k\in C$.
        Hence $k$ is a natural number and
            $\ringPower{R}{\ringNeg[R](\ringOne[R])}{k}
                \in\{\ringOne[R],\ringNeg[R](\ringOne[R])\}$.
        Hence $\ringPower{R}{\ringNeg[R](\ringOne[R])}{k}=\ringOne[R]$ or
            $\ringPower{R}{\ringNeg[R](\ringOne[R])}{k}
                =\ringNeg[R](\ringOne[R])$ by \auto{upair_iff}.
        $\ringPower{R}{\ringNeg[R](\ringOne[R])}{\suc{k}}
            =\ringMul[R](
                \ringPower{R}{\ringNeg[R](\ringOne[R])}{k},
                \ringNeg[R](\ringOne[R]))$ by
            \obv{ring_power_successor}.
        $\suc{k}$ is a natural number by \obv{omega_successor}.
        \begin{byCase}
            \caseOf{$\ringPower{R}{\ringNeg[R](\ringOne[R])}{k}
                    =\ringOne[R]$.}
                $\ringMul[R](\ringOne[R],\ringNeg[R](\ringOne[R]))
                    =\ringNeg[R](\ringOne[R])$ by
                    \obv{ring_mul_one_left}.
                Hence
                    $\ringPower{R}{\ringNeg[R](\ringOne[R])}{\suc{k}}
                        =\ringNeg[R](\ringOne[R])$.
                Follows by \auto{upair_iff}.
            \caseOf{$\ringPower{R}{\ringNeg[R](\ringOne[R])}{k}
                    \neq\ringOne[R]$.}
                Hence $\ringPower{R}{\ringNeg[R](\ringOne[R])}{k}
                    =\ringNeg[R](\ringOne[R])$.
                $\ringMul[R](\ringNeg[R](\ringOne[R]),
                        \ringNeg[R](\ringOne[R]))
                    =\ringNeg[R](\ringMul[R](\ringOne[R],
                        \ringNeg[R](\ringOne[R])))$ by
                    \obv{ring_neg_product_left}.
                $\ringMul[R](\ringOne[R],\ringNeg[R](\ringOne[R]))
                    =\ringNeg[R](\ringOne[R])$ by
                    \obv{ring_mul_one_left}.
                $\ringNeg[R](\ringNeg[R](\ringOne[R]))=\ringOne[R]$ by
                    \obv{ring_neg_involutive}.
                Hence
                    $\ringPower{R}{\ringNeg[R](\ringOne[R])}{\suc{k}}
                        =\ringOne[R]$.
                Follows by \obv{upair_iff}.
        \end{byCase}
    \end{subproof}
    Hence $C$ is an inductive set by \obv{inductive_set}.
    Hence every natural number is an element of $C$ by
        \obv{omega_minimal}.
    Hence $n\in C$.
    Follows by \auto{upair_iff}.
\end{proof}

\begin{corollary}\label{ring_power_of_negative_one_in_pair}
    Let $R$ be a ring and $n$ be a natural number.
    Then $\ringPower{R}{\ringNeg[R](\ringOne[R])}{n}
        \in\{\ringOne[R],\ringNeg[R](\ringOne[R])\}$.
\end{corollary}
\begin{proof}
    Follows by \auto{ring_power_of_negative_one,upair_iff}.
\end{proof}

\begin{proposition}\label{ring_power_of_one}
    Let $R$ be a ring.
    Let $n$ be a natural number.
    Then $\ringPower{R}{\ringOne[R]}{n}=\ringOne[R]$.
\end{proposition}
\begin{proof}
    Let $M=\MultiplicativeMonoid{R}$.
    $M$ is a monoid by
        \obv{multiplicative_monoid_is_exact_monoid}.
    $\neutral[M]=\ringOne[R]$ by
        \auto{multiplicative_monoid_fields}.
    $\monoidPower{M}{\neutral[M]}{n}=\neutral[M]$ by
        \obv{monoid_power_neutral}.
    Follows by \obv{ring_power}.
\end{proof}

\begin{proposition}\label{ring_power_of_zero_positive}
    Let $R$ be a ring.
    Let $n$ be a natural number.
    Suppose $n\neq\emptyset$.
    Then $\ringPower{R}{\ringZero[R]}{n}=\ringZero[R]$.
\end{proposition}
\begin{proof}
    Take a natural number $k$ such that $n=\suc{k}$ by
        \obv{natural_zero_or_successor}.
    $\ringZero[R]\in\carrier[R]$ by \obv{ring_zero_type}.
    $\ringPower{R}{\ringZero[R]}{k}\in\carrier[R]$ by
        \obv{ring_power_type}.
    $\ringPower{R}{\ringZero[R]}{\suc{k}}
        =\ringMul[R](\ringPower{R}{\ringZero[R]}{k},\ringZero[R])$ by
        \obv{ring_power_successor}.
    Hence $\ringPower{R}{\ringZero[R]}{\suc{k}}=\ringZero[R]$ by
        \obv{ring_mul_zero_right}.
\end{proof}

\begin{theorem}\label{ring_power_addition}
    Let $R$ be a ring.
    Let $x\in\carrier[R]$.
    Let $m,n$ be natural numbers.
    Then $\ringPower{R}{x}{m+n}=
        \ringMul[R](\ringPower{R}{x}{m},\ringPower{R}{x}{n})$.
\end{theorem}
\begin{proof}
    $\MultiplicativeMonoid{R}$ is a monoid by
        \obv{multiplicative_monoid_is_exact_monoid}.
    $x\in\carrier[\MultiplicativeMonoid{R}]$ by
        \auto{multiplicative_monoid_fields}.
    $\monoidPower{\MultiplicativeMonoid{R}}{x}{m+n}=
        \mul[\MultiplicativeMonoid{R}](
            \monoidPower{\MultiplicativeMonoid{R}}{x}{m},
            \monoidPower{\MultiplicativeMonoid{R}}{x}{n})$ by
        \obv{monoid_power_addition}.
    Follows by \auto{ring_power,multiplicative_monoid_fields}.
\end{proof}

\begin{theorem}\label{ring_power_multiplication}
    Let $R$ be a ring.
    Let $x\in\carrier[R]$.
    Let $m,n$ be natural numbers.
    Then $\ringPower{R}{x}{\naturalMul{m}{n}}=
        \ringPower{R}{\ringPower{R}{x}{m}}{n}$.
\end{theorem}
\begin{proof}
    $\MultiplicativeMonoid{R}$ is a monoid by
        \obv{multiplicative_monoid_is_exact_monoid}.
    $x\in\carrier[\MultiplicativeMonoid{R}]$ by
        \auto{multiplicative_monoid_fields}.
    $\monoidPower{\MultiplicativeMonoid{R}}{x}{
            \naturalMul{m}{n}}=
        \monoidPower{\MultiplicativeMonoid{R}}{
            \monoidPower{\MultiplicativeMonoid{R}}{x}{m}}{n}$ by
        \obv{monoid_power_multiplication}.
    Follows by \auto{ring_power}.
\end{proof}

\begin{theorem}\label{ring_power_of_commuting_product}
    Let $R$ be a ring and $x,y\in\carrier[R]$.
    Suppose $\ringMul[R](x,y)=\ringMul[R](y,x)$.
    Let $n$ be a natural number.
    Then
        $\ringPower{R}{\ringMul[R](x,y)}{n}
            =\ringMul[R](\ringPower{R}{x}{n},
                \ringPower{R}{y}{n})$.
\end{theorem}
\begin{proof}
    Let $M=\MultiplicativeMonoid{R}$.
    $M$ is a monoid by
        \obv{multiplicative_monoid_is_exact_monoid}.
    $\carrier[M]=\carrier[R]$ and $\mul[M]=\ringMul[R]$ by
        \auto{multiplicative_monoid_fields}.
    Hence $x,y\in\carrier[M]$ and
        $\mul[M](x,y)=\mul[M](y,x)$.
    $\monoidPower{M}{\mul[M](x,y)}{n}
        =\mul[M](\monoidPower{M}{x}{n},
            \monoidPower{M}{y}{n})$ by
        \obv{monoid_power_of_commuting_product}.
    Follows by
        \auto{ring_power,multiplicative_monoid_fields}.
\end{proof}

\begin{corollary}\label{commutative_ring_power_of_product}
    Let $R$ be a commutative ring and $x,y\in\carrier[R]$.
    Let $n$ be a natural number.
    Then
        $\ringPower{R}{\ringMul[R](x,y)}{n}
            =\ringMul[R](\ringPower{R}{x}{n},
                \ringPower{R}{y}{n})$.
\end{corollary}
\begin{proof}
    $R$ is a ring by \obv{commutative_ring_is_ring}.
    $\ringMul[R](x,y)=\ringMul[R](y,x)$ by
        \obv{commutative_ring_mul_commutative}.
    Follows by \obv{ring_power_of_commuting_product}.
\end{proof}

\begin{theorem}\label{commutative_ring_power_after_multiplying_inverse}
    Let $R$ be a commutative ring and $x,u,v,c\in\carrier[R]$.
    Suppose $\ringMul[R](u,v)=\ringOne[R]$.
    Let $n$ be a natural number.
    Suppose $\ringPower{R}{x}{n}
        =\ringMul[R](\ringPower{R}{u}{n},c)$.
    Then $\ringPower{R}{\ringMul[R](x,v)}{n}=c$.
\end{theorem}
\begin{proof}
    $R$ is a ring by \obv{commutative_ring_is_ring}.
    $\ringPower{R}{x}{n},\ringPower{R}{u}{n},
        \ringPower{R}{v}{n}\in\carrier[R]$ by
        \obv{ring_power_type}.
    $\ringPower{R}{\ringMul[R](x,v)}{n}
        =\ringMul[R](\ringPower{R}{x}{n},\ringPower{R}{v}{n})$ by
        \obv{commutative_ring_power_of_product}.
    $\ringPower{R}{\ringMul[R](u,v)}{n}
        =\ringMul[R](\ringPower{R}{u}{n},\ringPower{R}{v}{n})$ by
        \obv{commutative_ring_power_of_product}.
    $\ringPower{R}{\ringOne[R]}{n}=\ringOne[R]$ by
        \obv{ring_power_of_one}.
    Hence $\ringMul[R](\ringPower{R}{u}{n},\ringPower{R}{v}{n})
        =\ringOne[R]$.
    $\ringMul[R](\ringPower{R}{u}{n},c)
        \in\carrier[R]$ by \obv{ring_mul_type}.
    $\ringMul[R](\ringMul[R](\ringPower{R}{u}{n},c),
            \ringPower{R}{v}{n})
        =\ringMul[R](\ringPower{R}{u}{n},
            \ringMul[R](c,\ringPower{R}{v}{n}))$ by
        \obv{ring_mul_associative}.
    $\ringMul[R](c,\ringPower{R}{v}{n})
        =\ringMul[R](\ringPower{R}{v}{n},c)$ by
        \obv{commutative_ring_mul_commutative}.
    $\ringMul[R](\ringPower{R}{u}{n},
            \ringMul[R](\ringPower{R}{v}{n},c))
        =\ringMul[R](
            \ringMul[R](\ringPower{R}{u}{n},\ringPower{R}{v}{n}),c)$ by
        \obv{ring_mul_associative}.
    Hence $\ringMul[R](\ringMul[R](\ringPower{R}{u}{n},c),
            \ringPower{R}{v}{n})
        =\ringMul[R](\ringOne[R],c)$.
    Hence $\ringMul[R](\ringMul[R](\ringPower{R}{u}{n},c),
            \ringPower{R}{v}{n})=c$ by
        \obv{ring_mul_one_left}.
\end{proof}
